% Source: https://artofproblemsolving.com/community/c3h1840861_number_theory
% Stable identity: aops-topic-1840861-q01
\textbf{First inequality.} Each $a_i$ divides $L$ and both are positive, so $a_i\le L$. Adding these $n$ inequalities, $a_1+\cdots+a_n\le nL$, i.e.\ the mean is at most $L$. Equality forces $a_i=L$ for every $i$, so all the $a_i$ are equal (and conversely, if all are equal to $a$, then $L=a$ is the mean).

\textbf{Second inequality.} $G$ divides each $a_i$, so $G\le a_i$. Adding, $a_1+\cdots+a_n\ge nG$, i.e.\ the mean is at least $G$. Equality forces $a_i=G$ for every $i$, so all the $a_i$ are equal (and conversely). $\blacksquare$

(The proof in the post, which writes $L=\dfrac{a_1b_1+\cdots+a_nb_n}{n}$ with $b_i=\dfrac{L}{a_i}\ge1$, is the same argument and is correct.)
